% ============================================================
% FASE 2 + 3 DEL HITO 5: DETECCION DE DRIFT + ALERTAS + SIMULACION
% Bloque listo para integrar en la memoria del proyecto VinOps
% ============================================================

\subsection{Fase 2 -- Deteccion de Deriva de Datos (Data Drift)}

\subsubsection{Fundamento teorico}

El \textit{data drift} o deriva de datos se define como el cambio en la distribucion de las variables de entrada ($P(\mathbf{X})$) entre el entrenamiento y la operacion. A diferencia del \textit{concept drift} ---que afecta la relacion $P(Y|\mathbf{X})$---, el data drift es detectable sin \textit{ground truth}, lo cual lo hace especialmente relevante para VinOps, donde la etiqueta real del cultivar no esta disponible en el momento de la prediccion.

Se distinguen tres niveles de severidad, siguiendo la taxonomia de \cite{alibi_detect} y \cite{nannyml}:

\begin{itemize}
    \item \textbf{Drift leve} ($0.1 < PSI < 0.25$): monitorizar, no requiere accion inmediata.
    \item \textbf{Drift moderado} ($PSI > 0.25$ o $KS\_p < 0.05$ en 1--2 variables): revisar proceso de produccion.
    \item \textbf{Drift severo} ($KS\_p < 0.05$ en $\geq 3$ variables o $W > 2.0$): reentrenamiento recomendado.
\end{itemize}

\subsubsection{Tecnicas implementadas}

Se han implementado tres tecnicas complementarias, todas basadas en \texttt{scipy} para mantener el contenedor ligero:

\begin{table}[htbp]
\centering
\caption{Tecnicas de deteccion de drift implementadas en VinOps}
\label{tab:tecnicas-drift}
\begin{tabular}{@{}llp{5.5cm}@{}}
\toprule
\textbf{Tecnica} & \textbf{Tipo} & \textbf{Umbrales} \\
\midrule
Kolmogorov-Smirnov (KS) & Univariado & $p < 0.05$ por variable \\
Population Stability Index (PSI) & Univariado & $PSI > 0.25$ (alerta), $> 0.1$ (monitor) \\
Wasserstein (media estandarizada) & Multivariado & $W > 2.0$ (alerta), $> 1.0$ (monitor) \\
\bottomrule
\end{tabular}
\end{table}

\paragraph{KS test.} Contrasta la hipotesis nula de que dos muestras provienen de la misma distribucion continua. Se aplica por variable entre el dataset de referencia (entrenamiento) y la ventana deslizante de produccion.

\paragraph{PSI.} Mide la magnitud del desplazamiento de la distribucion dividiendo en 10 bins y comparando proporciones. Es la metrica estandar en modelos de scoring y riesgo crediticio.

\paragraph{Wasserstein multivariado.} Calcula la distancia de Wasserstein por variable, normalizada por la desviacion estandar de referencia, y promedia. Detecta cambios conjuntos que las tecnicas univariadas podrian omitir.

\subsubsection{Ventana deslizante y evaluacion periodica}

Dado el volumen reducido del dataset (178 instancias), se configura una ventana de 20 muestras con evaluacion cada 10 predicciones. Esto equilibra:

\begin{itemize}
    \item \textbf{Sensibilidad:} 20 muestras son suficientes para estimar 13 variables.
    \item \textbf{Latencia:} evaluar cada 10 predicciones evita sobrecarga computacional.
    \item \textbf{Memoria:} la ventana se mantiene en memoria (deque) sin persistencia externa.
\end{itemize}

\subsubsection{Endpoint /drift}

Devuelve el estado completo del detector:

\begin{itemize}
    \item \texttt{status}: OK / MONITOR / DRIFT\_DETECTED
    \item \texttt{recommendation}: OK / MONITOR / RETRAIN
    \item \texttt{ks\_test}: p-values por variable, variables alertadas
    \item \texttt{psi}: scores por variable, variables alertadas
    \item \texttt{wasserstein}: distancia multivariada
\end{itemize}

\subsection{Fase 3 -- Sistema de Alertas}

\subsubsection{Canal de alertas: Telegram Bot}

Se ha seleccionado Telegram como canal de notificaciones por las siguientes razones:

\begin{enumerate}
    \item \textbf{Simplicidad tecnica:} la API de Telegram es HTTP REST sin autenticacion compleja (solo token).
    \item \textbf{Visibilidad:} el grupo incluye al equipo tecnico y al stakeholder (director tecnico de la bodega).
    \item \textbf{Demo en vivo:} durante la presentacion, el tribunal puede ver las alertas en tiempo real en el chat.
    \item \textbf{Sin infraestructura propia:} no requiere servidor SMTP, conectores de Teams ni permisos de tenant.
\end{enumerate}

\subsubsection{Throttling}

Para evitar spam en escenarios de degradacion continua, se implementa un mecanismo de \textit{throttling} de 60 segundos por tipo de alerta. Esto garantiza que el canal no se sature si el sistema detecta drift persistente.

\subsubsection{Tipos de alerta}

\begin{table}[htbp]
\centering
\caption{Tipos de alerta implementados en VinOps}
\label{tab:alertas}
\begin{tabular}{@{}llp{5cm}l@{}}
\toprule
\textbf{Tipo} & \textbf{Trigger} & \textbf{Mensaje} & \textbf{Canal} \\
\midrule
Drift & $KS\_p < 0.05$ en $\geq 3$ vars o $PSI > 0.25$ en $\geq 2$ vars & Deriva detectada, reentrenar & Telegram \\
Anomalia & Tasa de anomalias $> 25\%$ en ventana & Tasa de anomalias elevada & Telegram \\
Confianza baja & Confianza media $< 0.65$ en ventana & Confianza media baja & Telegram \\
Operativa & CPU $> 80\%$ o RAM $> 300$ MB & Alerta de recurso & Telegram \\
Simulacion & Endpoint /simulate/* invocado & Simulacion completada & Telegram \\
\bottomrule
\end{tabular}
\end{table}

\subsubsection{Simulacion de entorno anomalo}

Se han implementado dos endpoints de simulacion para validar el sistema sin datos reales de produccion:

\paragraph{/simulate/anomaly.} Genera $N$ muestras con valores extremos ($\pm 2.5\sigma$ respecto a la media de referencia) pero dentro de los rangos fisicos del vino. Permite verificar que el detector de anomalias y el sistema de alertas responden correctamente.

\paragraph{/simulate/drift.} Genera $N$ muestras donde una variable especifica esta desplazada $+k\sigma$ respecto a la distribucion de entrenamiento, manteniendo el resto de variables normales. Simula un cambio real en la linea de produccion (p. ej. uva de otra zona con mayor contenido alcoholico).

Ambas simulaciones inyectan las muestras en la ventana de drift, evaluan el estado y envian una alerta resumen al canal de Telegram.

\subsubsection{Decisiones de diseno}

\begin{enumerate}
    \item \textbf{No se usa Alibi Detect ni NannyML:} aunque son herramientas de referencia del estado del arte \cite{alibi_detect, nannyml}, se opto por implementacion propia con \texttt{scipy} para mantener el contenedor ligero y evitar dependencias pesadas en un entorno academico. En produccion real se migraria a dichas herramientas.
    \item \textbf{Token en .env, no en codigo:} la seguridad del bot se gestiona via variables de entorno, con \texttt{.env} en \texttt{.gitignore}.
    \item \textbf{Metricas operativas y de modelo separadas:} siguiendo la taxonomia del profesor, las alertas operativas (CPU/RAM) y las de modelo (drift/confianza) tienen tipos distintos y umbrales independientes.
\end{enumerate}
